\begin{proof}

\begin{enumerate}[label={$\langle1\rangle$\arabic*.},leftmargin=*]
\item Let $D$ be the conductor divisor and
$C=\operatorname{Sing}(S)_{\mathrm{red}}$. By P4.12,
$K_T+D\equiv0$. Contract $T$ to a minimal rational surface
$T_0$ and push $D$ forward to $D_0$. At every reversed
blowup the exceptional coefficient is $\mu-1$, where
$\mu$ is the multiplicity of the current conductor at the
centre. In particular every centre lies on the conductor.

\item We give the complete minimal-surface case split. On
$\mathbb P^2$, the divisor $D_0$ is a plane cubic. An
irreducible cubic is smooth elliptic, rational nodal, or
rational cuspidal. A reducible cubic is a smooth conic
and a line, three lines, a double line and another line,
or a triple line. The conic and line meet at two
transverse points or at one tangency; three distinct
lines form a triangle or are concurrent. Thus either
$D_0$ is reduced nodal and has fundamental group
$\mathbb Z$ or $\mathbb Z^2$, or its reduced support is
simply connected and all normalized components are
rational.

\item On a Hirzebruch surface, write $F$ for a ruling fibre.
Since $D_0\cdot F=2$, its horizontal part is a double
section, two distinct reduced sections, or a reduced
irreducible bisection. A double section together with
vertical fibres is a rational tree. For the other cases
write
\[
 D_0=H_1+H_2+\sum_j a_jF_j
 \quad\hbox{or}\quad
 D_0=H+\sum_j a_jF_j.
\]
In the first case adjunction on either section gives
\[
 H_1\cdot H_2+\sum_j a_j=2.
\]
The support is therefore a rational tree or a single
rational cycle. A cycle occurs only with all coefficients
one and ordinary double intersection points. This
includes the cases of two section-intersection points,
one section-intersection point and one vertical fibre,
or two vertical fibres. Coincident intersections or a
vertical fibre through the section-intersection point
instead give trees.

\item For a bisection, adjunction gives
\[
 p_a(H)=1-\sum_j a_j.
\]
There is either no vertical component and $H$ has
arithmetic genus one, or one coefficient-one vertical
fibre and a smooth rational bisection. In the former
case $H$ is smooth elliptic, rational nodal, or rational
cuspidal. In the latter it meets the fibre at two distinct
points or at one tangency. Again the only nontrivial
fundamental groups are $\mathbb Z$ and $\mathbb Z^2$,
and those cases have reduced nodal conductor.

\item We follow these alternatives through the blowups. At a
point of conductor multiplicity at least two, the new
conductor contains the exceptional $\mathbb P^1$.
Its contraction is the quotient which collapses that
connected simply connected subcomplex to the old
point. Van Kampen therefore makes the induced map on
fundamental groups an isomorphism. At a point of
multiplicity one the conductor is smooth and reduced,
and its strict transform is isomorphic to it. Thus
$\pi_1(D_{\mathrm{red}})$ is unchanged throughout.
If the initial group was nontrivial, the conductor was
reduced nodal. A smooth-point blowup adds no conductor
component, and a node blowup inserts one component of
coefficient one. It remains reduced nodal.

\item In this latter case the generic transverse conductor
exponent of every normalization branch is one. The
exponent is its intrinsic conductor exponent plus the
intersection multiplicities with the other branches.
A singular branch has intrinsic exponent at least two.
Exponent one therefore means two smooth branches
meeting transversely. Consequently $S$ has ordinary
double curves outside finitely many points. P4.20 makes
$S$ normal crossing everywhere, and P4.6--P4.8 give
$e(S)=2$ for this rational normalization.

\item In the remaining case $D_{\mathrm{red}}$ is connected
and simply connected, and all its normalized components
are rational. The conductor pushout verified in P4R.3
gives
\[
 \pi_1(S)=\pi_1(T)*_{\pi_1(D_{\mathrm{red}})}\pi_1(C)
          =\pi_1(C).
\]
Every component of $C$ has rational normalization,
because it is the finite nonconstant image of a
rational component of $D_{\mathrm{red}}$. P4R.1 makes
$\pi_1(C)$ virtually abelian. For a connected curve
with rational normalized components, this is the free
group $F_b$ of its point-identification graph, so
$b\le1$. If $c\ge1$ is its number of components, then
\[
 e(C)=c+1-b\ge c\ge1.
\]
P4.13 gives $e(S_{\mathrm{reg}})\ge0$. Additivity proves
$e(S)=e(S_{\mathrm{reg}})+e(C)>0$.

\end{enumerate}
\end{proof}
